\documentclass[10pt,a4paper]{amsart}
\usepackage{amsmath,amssymb,amsthm,graphicx}
\usepackage{verbatim}
\newenvironment{reference}{\comment}{\endcomment}
%\newenvironment{reference}{}{}
\newenvironment{slogan}{\comment}{\endcomment}
\newenvironment{history}{\comment}{\endcomment}

% For commutative diagrams we use Xy-pic
\usepackage[all]{xy}

% We use 2cell for 2-commutative diagrams.
\xyoption{2cell}
% Theorem environments.
%
\theoremstyle{plain}
\newtheorem{theorem}[subsection]{Theorem}
\newtheorem{proposition}[subsection]{Proposition}
\newtheorem{lemma}[subsection]{Lemma}

\theoremstyle{definition}
\newtheorem{definition}[subsection]{Definition}
\newtheorem{example}[subsection]{Example}
\newtheorem{exercise}[subsection]{Exercise}
\newtheorem{situation}[subsection]{Situation}

\theoremstyle{remark}
\newtheorem{remark}[subsection]{Remark}
\newtheorem{remarks}[subsection]{Remarks}

\numberwithin{equation}{subsection}

% Macros
%
\def\lim{\mathop{\mathrm{lim}}\nolimits}
\def\colim{\mathop{\mathrm{colim}}\nolimits}
\def\Spec{\mathop{\mathrm{Spec}}}
\def\Hom{\mathop{\mathrm{Hom}}\nolimits}
\def\Ext{\mathop{\mathrm{Ext}}\nolimits}
\def\SheafHom{\mathop{\mathcal{H}\!\mathit{om}}\nolimits}
\def\SheafExt{\mathop{\mathcal{E}\!\mathit{xt}}\nolimits}
\def\Sch{\mathit{Sch}}
\def\Mor{\mathop{\mathrm{Mor}}\nolimits}
\def\Ob{\mathop{\mathrm{Ob}}\nolimits}
\def\Sh{\mathop{\mathit{Sh}}\nolimits}
\def\NL{\mathop{N\!L}\nolimits}
\def\CH{\mathop{\mathrm{CH}}\nolimits}
\def\proetale{{pro\text{-}\acute{e}tale}}
\def\etale{{\acute{e}tale}}
\def\QCoh{\mathit{QCoh}}
\def\Ker{\mathop{\mathrm{Ker}}}
\def\Im{\mathop{\mathrm{Im}}}
\def\Coker{\mathop{\mathrm{Coker}}}
\def\Coim{\mathop{\mathrm{Coim}}}

% Boxtimes
%
\DeclareMathSymbol{\boxtimes}{\mathbin}{AMSa}{"02}

%
% Macros for moduli stacks/spaces
%
\def\QCohstack{\mathcal{QC}\!\mathit{oh}}
\def\Cohstack{\mathcal{C}\!\mathit{oh}}
\def\Spacesstack{\mathcal{S}\!\mathit{paces}}
\def\Quotfunctor{\mathrm{Quot}}
\def\Hilbfunctor{\mathrm{Hilb}}
\def\Curvesstack{\mathcal{C}\!\mathit{urves}}
\def\Polarizedstack{\mathcal{P}\!\mathit{olarized}}
\def\Complexesstack{\mathcal{C}\!\mathit{omplexes}}
% \Pic is the operator that assigns to X its picard group, usage \Pic(X)
% \Picardstack_{X/B} denotes the Picard stack of X over B
% \Picardfunctor_{X/B} denotes the Picard functor of X over B
\def\Pic{\mathop{\mathrm{Pic}}\nolimits}
\def\Picardstack{\mathcal{P}\!\mathit{ic}}
\def\Picardfunctor{\mathrm{Pic}}
\def\Deformationcategory{\mathcal{D}\!\mathit{ef}}
\usepackage{cleveref}
\crefname{originaldefinition}{original definition}{original definitions}
\Crefname{originaldefinition}{Original definition}{Original definitions}
\makeatletter
\newlabel{spaces-morphisms-definition-scheme-theoretic-intersection-union}{{0CYZ}{source}{}{sourcepointer-1}{}}
\newlabel{spaces-morphisms-definition-scheme-theoretic-intersection-union@cref}{{[originaldefinition][0][]0CYZ}{[1][0][]source}{}{}{}}
\makeatother
\crefname{originallemma}{original lemma}{original lemmas}
\Crefname{originallemma}{Original lemma}{Original lemmas}
\makeatletter
\newlabel{spaces-morphisms-lemma-closed-immersion-proper}{{04XZ}{source}{}{sourcepointer-2}{}}
\newlabel{spaces-morphisms-lemma-closed-immersion-proper@cref}{{[originallemma][0][]04XZ}{[1][0][]source}{}{}{}}
\makeatother
\crefname{originallemma}{original lemma}{original lemmas}
\Crefname{originallemma}{Original lemma}{Original lemmas}
\makeatletter
\newlabel{spaces-morphisms-lemma-composition-finite-type}{{03XG}{source}{}{sourcepointer-3}{}}
\newlabel{spaces-morphisms-lemma-composition-finite-type@cref}{{[originallemma][0][]03XG}{[1][0][]source}{}{}{}}
\makeatother
\crefname{originallemma}{original lemma}{original lemmas}
\Crefname{originallemma}{Original lemma}{Original lemmas}
\makeatletter
\newlabel{spaces-morphisms-lemma-composition-proper}{{04XY}{source}{}{sourcepointer-4}{}}
\newlabel{spaces-morphisms-lemma-composition-proper@cref}{{[originallemma][0][]04XY}{[1][0][]source}{}{}{}}
\makeatother
\crefname{originallemma}{original lemma}{original lemmas}
\Crefname{originallemma}{Original lemma}{Original lemmas}
\makeatletter
\newlabel{spaces-morphisms-lemma-image-proper-is-proper}{{08AJ}{source}{}{sourcepointer-5}{}}
\newlabel{spaces-morphisms-lemma-image-proper-is-proper@cref}{{[originallemma][0][]08AJ}{[1][0][]source}{}{}{}}
\makeatother
\crefname{originallemma}{original lemma}{original lemmas}
\Crefname{originallemma}{Original lemma}{Original lemmas}
\makeatletter
\newlabel{spaces-morphisms-lemma-image-universally-closed-separated}{{05Z2}{source}{}{sourcepointer-6}{}}
\newlabel{spaces-morphisms-lemma-image-universally-closed-separated@cref}{{[originallemma][0][]05Z2}{[1][0][]source}{}{}{}}
\makeatother
\crefname{originallemma}{original lemma}{original lemmas}
\Crefname{originallemma}{Original lemma}{Original lemmas}
\makeatletter
\newlabel{spaces-morphisms-lemma-immersion-locally-finite-type}{{06ED}{source}{}{sourcepointer-7}{}}
\newlabel{spaces-morphisms-lemma-immersion-locally-finite-type@cref}{{[originallemma][0][]06ED}{[1][0][]source}{}{}{}}
\makeatother
\crefname{originallemma}{original lemma}{original lemmas}
\Crefname{originallemma}{Original lemma}{Original lemmas}
\makeatletter
\newlabel{spaces-morphisms-lemma-scheme-theoretic-union}{{0CZ1}{source}{}{sourcepointer-8}{}}
\newlabel{spaces-morphisms-lemma-scheme-theoretic-union@cref}{{[originallemma][0][]0CZ1}{[1][0][]source}{}{}{}}
\makeatother
\crefname{originaldefinition}{original definition}{original definitions}
\Crefname{originaldefinition}{Original definition}{Original definitions}
\makeatletter
\newlabel{spaces-properties-definition-reduced-induced-space}{{047X}{source}{}{sourcepointer-9}{}}
\newlabel{spaces-properties-definition-reduced-induced-space@cref}{{[originaldefinition][0][]047X}{[1][0][]source}{}{}{}}
\makeatother
\title{Properness over the base of a closed subset is independent of the closed subspace structure}
\author{The Stacks Project Authors}
\date{stacks-project:spaces-perfect.tex:local-snapshot-20261012; GNU FDL 1.2 or later}
\begin{document}
\maketitle
\noindent\textit{Editorial note.} All delivered mathematical text is the author's own, copied from the pinned Stacks Project chapter file spaces-perfect.tex; no statement or proof step was rewritten, reordered or removed.\par\medskip
\noindent\textit{Editorial note.} The counted proof is the author's complete argument: the proof is complete: the three implications are all carried out, with the union trick doing the work\par\medskip
\noindent\textit{Editorial note.} The label lemma-closed-proper-over-base names the counted statement, and the proof environment immediately following it in the source is delivered verbatim as its proof.\par\medskip
\noindent\textit{Editorial note.} The counted statement is carried by the author's own lemma environment in the chapter source.\par\medskip
\noindent\textit{Allowed standard prerequisites (editorial).}\begin{itemize}
\item reduced induced closed subspace structures and scheme-theoretic unions of closed subspaces
\item properness, separatedness and quasi-compactness of morphisms of algebraic spaces and their permanence under closed immersions, compositions and images of universally closed maps
\item bijective closed immersions as universal homeomorphisms
\end{itemize}\medskip
\ifdefined\hypertarget\hypertarget{sourcepointer-1}{}\fi
\noindent\textit{Source reference.} Original definition 0CYZ: sources/stacks-chow/tags.txt, source lines 16259--16259. Prerequisite pointer invoked at proof line(s) 1057; it names a standard fact of a sibling chapter that the counted argument uses, and every step of the counted proof built on it is retained verbatim in the delivered body.\par\medskip
\ifdefined\hypertarget\hypertarget{sourcepointer-2}{}\fi
\noindent\textit{Source reference.} Original lemma 04XZ: sources/stacks-chow/tags.txt, source lines 6412--6412. Prerequisite pointer invoked at proof line(s) 1064; it names a standard fact of a sibling chapter that the counted argument uses, and every step of the counted proof built on it is retained verbatim in the delivered body.\par\medskip
\ifdefined\hypertarget\hypertarget{sourcepointer-3}{}\fi
\noindent\textit{Source reference.} Original lemma 03XG: sources/stacks-chow/tags.txt, source lines 5147--5147. Prerequisite pointer invoked at proof line(s) 1076; it names a standard fact of a sibling chapter that the counted argument uses, and every step of the counted proof built on it is retained verbatim in the delivered body.\par\medskip
\ifdefined\hypertarget\hypertarget{sourcepointer-4}{}\fi
\noindent\textit{Source reference.} Original lemma 04XY: sources/stacks-chow/tags.txt, source lines 6411--6411. Prerequisite pointer invoked at proof line(s) 1065; it names a standard fact of a sibling chapter that the counted argument uses, and every step of the counted proof built on it is retained verbatim in the delivered body.\par\medskip
\ifdefined\hypertarget\hypertarget{sourcepointer-5}{}\fi
\noindent\textit{Source reference.} Original lemma 08AJ: sources/stacks-chow/tags.txt, source lines 10534--10534. Prerequisite pointer invoked at proof line(s) 1081; it names a standard fact of a sibling chapter that the counted argument uses, and every step of the counted proof built on it is retained verbatim in the delivered body.\par\medskip
\ifdefined\hypertarget\hypertarget{sourcepointer-6}{}\fi
\noindent\textit{Source reference.} Original lemma 05Z2: sources/stacks-chow/tags.txt, source lines 7676--7676. Prerequisite pointer invoked at proof line(s) 1071; it names a standard fact of a sibling chapter that the counted argument uses, and every step of the counted proof built on it is retained verbatim in the delivered body.\par\medskip
\ifdefined\hypertarget\hypertarget{sourcepointer-7}{}\fi
\noindent\textit{Source reference.} Original lemma 06ED: sources/stacks-chow/tags.txt, source lines 8212--8212. Prerequisite pointer invoked at proof line(s) 1075; it names a standard fact of a sibling chapter that the counted argument uses, and every step of the counted proof built on it is retained verbatim in the delivered body.\par\medskip
\ifdefined\hypertarget\hypertarget{sourcepointer-8}{}\fi
\noindent\textit{Source reference.} Original lemma 0CZ1: sources/stacks-chow/tags.txt, source lines 16261--16261. Prerequisite pointer invoked at proof line(s) 1060; it names a standard fact of a sibling chapter that the counted argument uses, and every step of the counted proof built on it is retained verbatim in the delivered body.\par\medskip
\ifdefined\hypertarget\hypertarget{sourcepointer-9}{}\fi
\noindent\textit{Source reference.} Original definition 047X: sources/stacks-chow/tags.txt, source lines 5524--5524. Prerequisite pointer invoked at proof line(s) 1040; it names a standard fact of a sibling chapter that the counted argument uses, and every step of the counted proof built on it is retained verbatim in the delivered body.\par\medskip
\begin{lemma}
\label{lemma-closed-proper-over-base}
Let $S$ be a scheme. Let $f : X \to Y$ be a morphism of algebraic spaces
over $S$ which is locally of finite type. Let $T \subset |X|$ be a closed
subset. The following are equivalent
\begin{enumerate}
\item the morphism $Z \to Y$ is proper if $Z$ is the reduced
induced algebraic space structure on $T$
(Properties of Spaces, Definition
\ref{spaces-properties-definition-reduced-induced-space}),
\item for some closed subspace $Z \subset X$ with $|Z| = T$
the morphism $Z \to Y$ is proper, and
\item for any closed subspace $Z \subset X$ with $|Z| = T$ the morphism
$Z \to Y$ is proper.
\end{enumerate}
\end{lemma}

\begin{proof}
The implications (3) $\Rightarrow$ (1) and (1) $\Rightarrow$ (2)
are immediate. Thus it suffices to prove that (2) implies (3).
We urge the reader to find their own proof of this fact.
Let $Z'$ and $Z''$ be closed subspaces with $T = |Z'| = |Z''|$
such that $Z' \to Y$ is a proper morphism of algebraic spaces.
We have to show that $Z'' \to Y$ is proper too.
Let $Z''' = Z' \cup Z''$ be the scheme theoretic union, see
Morphisms of Spaces, Definition
\ref{spaces-morphisms-definition-scheme-theoretic-intersection-union}.
Then $Z'''$ is another closed subspace with $|Z'''| = T$.
This follows for example from the description of scheme theoretic unions in
Morphisms of Spaces, Lemma \ref{spaces-morphisms-lemma-scheme-theoretic-union}.
Since $Z'' \to Z'''$ is a closed immersion it suffices to prove
that $Z''' \to Y$ is proper (see
Morphisms of Spaces, Lemmas
\ref{spaces-morphisms-lemma-closed-immersion-proper} and
\ref{spaces-morphisms-lemma-composition-proper}).
The morphism $Z' \to Z'''$ is a bijective closed immersion
and in particular surjective and universally closed.
Then the fact that $Z' \to Y$ is separated implies that
$Z''' \to Y$ is separated, see
Morphisms of Spaces, Lemma
\ref{spaces-morphisms-lemma-image-universally-closed-separated}.
Moreover $Z''' \to Y$ is locally of finite type
as $X \to Y$ is locally of finite type
(Morphisms of Spaces, Lemmas
\ref{spaces-morphisms-lemma-immersion-locally-finite-type} and
\ref{spaces-morphisms-lemma-composition-finite-type}).
Since $Z' \to Y$ is quasi-compact and $Z' \to Z'''$ is a
universal homeomorphism we see that $Z''' \to Y$ is quasi-compact.
Finally, since $Z' \to Y$ is universally closed, we see that
the same thing is true for $Z''' \to Y$ by
Morphisms of Spaces, Lemma \ref{spaces-morphisms-lemma-image-proper-is-proper}.
This finishes the proof.
\end{proof}

\end{document}
